% Canzoniere di Clan - template. I segnaposto tra doppi angoli (maiuscoli)
% sono sostituiti da scripts/generate_latex.py. Compilare con pdflatex (due passate) o latexmk.
\documentclass[a5paper,landscape,10pt,twoside]{article}

\usepackage[utf8]{inputenc}
\usepackage[T1]{fontenc}
\usepackage[italian]{babel}
\usepackage{lmodern}
\usepackage[a5paper,landscape,twoside,
  left=15mm,right=15mm,top=12mm,bottom=12mm,bindingoffset=15mm,
  includefoot,nohead,footskip=6mm]{geometry}
\usepackage{multicol}
\usepackage{graphicx}
\usepackage{xcolor}
\usepackage{url}
\usepackage{microtype}

\setlength{\parindent}{0pt}
\setlength{\parskip}{0pt}
\setlength{\columnsep}{8mm}
\raggedbottom
\sloppy
\hyphenpenalty=1000

\definecolor{accordo}{gray}{0.25}

% Interlinea ampia: lascia spazio agli accordi sopra il testo.
\makeatletter
\renewcommand{\normalsize}{\@setfontsize\normalsize{10}{17}}

% Stile pagina: solo numero in basso, sul lato esterno.
\def\ps@canz{%
  \def\@oddhead{}\def\@evenhead{}%
  \def\@oddfoot{\hfill\small\thepage}%
  \def\@evenfoot{\small\thepage\hfill}}

% \cantotesta{id}{titolo}{categoria}{autore}: nuova pagina + etichetta per gli indici
\newcommand{\cantotesta}[4]{%
  \clearpage
  \def\@currentlabel{#1}\label{canto:#1}%
  \noindent\hangindent=14mm\hangafter=1
  \hbox to 14mm{\fontsize{22}{24}\selectfont\bfseries #1\hss}%
  {\fontsize{15}{18}\selectfont\bfseries #2\par}
  {\small\scshape #3\ifx\relax#4\relax\else\ \textperiodcentered\ \textit{#4}\fi\par}
  \vspace{1mm}\hrule height 0.4pt\vspace{3mm}}
\makeatother
\pagestyle{canz}

% Riga di testo con accordi: l'accordo sta sopra la sillaba che lo segue.
\newcommand{\ch}[2]{\leavevmode\vbox{\offinterlineskip
  \hbox{\fontsize{8}{9}\selectfont\bfseries\color{accordo}#1\kern.3em}%
  \kern.45ex
  \hbox{#2}}}
\newcommand{\rigo}[1]{{\hangindent=3mm\hangafter=1\noindent #1\par}}
\newcommand{\strofa}{\par\vspace{2.4mm}}
\newcommand{\commento}[1]{{\par\small\itshape\color{gray}#1\par\vspace{1mm}}}
\newenvironment{ritornello}{\par\begingroup\itshape\leftskip=3mm\relax}{\par\endgroup}

% Indici
\newcommand{\titolosezione}[1]{\clearpage
  {\noindent\fontsize{16}{18}\selectfont\bfseries #1\par}
  \vspace{1mm}\hrule height 0.4pt\vspace{3mm}}
\newcommand{\lettera}[1]{\par\vspace{1.5mm}{\bfseries\color{accordo}#1}\par\nopagebreak}
\newcommand{\categoriaindice}[1]{\par\vspace{2mm}{\scshape\bfseries #1}\par\nopagebreak}
\newcommand{\voce}[2]{\noindent\hangindent=2.4em\hangafter=1
  \hbox to 2em{\hfil\bfseries #1}\hskip .4em #2\nobreak
  \leaders\hbox to .7em{\hss.\hss}\hfill\nobreak
  \hbox to 2em{\hfill\pageref{canto:#1}}\par}

\begin{document}

% ---- Frontespizio (numerazione romana, nessun numero stampato)
\pagenumbering{roman}
\thispagestyle{empty}
\begin{center}
  \vspace*{\fill}
  {\fontsize{40}{44}\selectfont\bfseries <<TITOLO>>\par}
  \vspace{4mm}
  {\Large <<SOTTOTITOLO>>\par}
  \vspace{10mm}
  {\large\itshape <<DEDICA>>\par}
  \vspace*{\fill}
  {\small <<ANNO>>\par}
\end{center}
<<NOTA>>

% ---- Canti (la numerazione araba riparte da 1: pagina stampata = campo "pagina")
\pagenumbering{arabic}
<<CANTI>>

% ---- Indici
<<INDICE_ALFABETICO>>
<<INDICE_TEMATICO>>

% ---- Ultima pagina: QR code verso la web app
\clearpage
\thispagestyle{empty}
\noindent
\begin{minipage}[c]{0.38\textwidth}
  \centering
  \includegraphics[width=44mm]{<<QR_FILE>>}
\end{minipage}\hfill
\begin{minipage}[c]{0.58\textwidth}
  {\fontsize{16}{19}\selectfont\bfseries Ascolta i canti e scarica gli spartiti\par}
  \vspace{3mm}
  Inquadra il codice con la fotocamera del telefono: si apre il canzoniere
  online, con l'audio di ogni canto, gli spartiti in PDF e la modalit\`a notte
  per il bivacco.\par
  \vspace{4mm}
  {\small\textbf{In caso di problemi con il QR\ldots}\par
  \begin{enumerate}\itemsep=0pt\topsep=1mm
    \item apri il browser del telefono (Chrome, Safari\ldots);
    \item scrivi nella barra degli indirizzi:\\
          {\small\url{<<URL>>}}
    \item cerca il canto per numero (quello stampato accanto al titolo);
    \item per riaprirlo senza connessione, aggiungilo alla schermata Home.
  \end{enumerate}}
\end{minipage}

\end{document}
